%Version 3.1 December 2024
% See section 11 of the User Manual for version history
%
%%%%%%%%%%%%%%%%%%%%%%%%%%%%%%%%%%%%%%%%%%%%%%%%%%%%%%%%%%%%%%%%%%%%%%
%%                                                                 %%
%% Please do not use \input{...} to include other tex files.       %%
%% Submit your LaTeX manuscript as one .tex document.              %%
%%                                                                 %%
%% All additional figures and files should be attached             %%
%% separately and not embedded in the \TeX\ document itself.       %%
%%                                                                 %%
%%%%%%%%%%%%%%%%%%%%%%%%%%%%%%%%%%%%%%%%%%%%%%%%%%%%%%%%%%%%%%%%%%%%%

%%\documentclass[referee,sn-basic]{sn-jnl}% referee option is meant for double line spacing

%%=======================================================%%
%% to print line numbers in the margin use lineno option %%
%%=======================================================%%

%%\documentclass[lineno,pdflatex,sn-basic]{sn-jnl}% Basic Springer Nature Reference Style/Chemistry Reference Style

%%=========================================================================================%%
%% the documentclass is set to pdflatex as default. You can delete it if not appropriate.  %%
%%=========================================================================================%%

%%\documentclass[sn-basic]{sn-jnl}% Basic Springer Nature Reference Style/Chemistry Reference Style

%%Note: the following reference styles support Namedate and Numbered referencing. By default the style follows the most common style. To switch between the options you can add or remove “Numbered” in the optional parenthesis. 
%%The option is available for: sn-basic.bst, sn-chicago.bst%  
 
%%\documentclass[pdflatex,sn-nature]{sn-jnl}% Style for submissions to Nature Portfolio journals
%%\documentclass[pdflatex,sn-basic]{sn-jnl}% Basic Springer Nature Reference Style/Chemistry Reference Style
\documentclass[pdflatex,sn-mathphys-num]{sn-jnl}% Math and Physical Sciences Numbered Reference Style
%%\documentclass[pdflatex,sn-mathphys-ay]{sn-jnl}% Math and Physical Sciences Author Year Reference Style
%%\documentclass[pdflatex,sn-aps]{sn-jnl}% American Physical Society (APS) Reference Style
%%\documentclass[pdflatex,sn-vancouver-num]{sn-jnl}% Vancouver Numbered Reference Style
%%\documentclass[pdflatex,sn-vancouver-ay]{sn-jnl}% Vancouver Author Year Reference Style
%%\documentclass[pdflatex,sn-apa]{sn-jnl}% APA Reference Style
%%\documentclass[pdflatex,sn-chicago]{sn-jnl}% Chicago-based Humanities Reference Style

%%%% Standard Packages
%%<additional latex packages if required can be included here>

\usepackage{graphicx}%
\usepackage{multirow}%
\usepackage{amsmath,amssymb,amsfonts}%
\usepackage{amsthm}%
\usepackage{mathrsfs}%
\usepackage[title]{appendix}%
\usepackage{xcolor}%
\usepackage{textcomp}%
\usepackage{manyfoot}%
\usepackage{booktabs}%
\usepackage{algorithm}%
\usepackage{algorithmicx}%
\usepackage{algpseudocode}%
\usepackage{listings}%
%%%%

%%%%%=============================================================================%%%%
%%%%  Remarks: This template is provided to aid authors with the preparation
%%%%  of original research articles intended for submission to journals published 
%%%%  by Springer Nature. The guidance has been prepared in partnership with 
%%%%  production teams to conform to Springer Nature technical requirements. 
%%%%  Editorial and presentation requirements differ among journal portfolios and 
%%%%  research disciplines. You may find sections in this template are irrelevant 
%%%%  to your work and are empowered to omit any such section if allowed by the 
%%%%  journal you intend to submit to. The submission guidelines and policies 
%%%%  of the journal take precedence. A detailed User Manual is available in the 
%%%%  template package for technical guidance.
%%%%%=============================================================================%%%%

%% as per the requirement new theorem styles can be included as shown below
\theoremstyle{thmstyleone}%
\newtheorem{theorem}{Theorem}%  meant for continuous numbers
%%\newtheorem{theorem}{Theorem}[section]% meant for sectionwise numbers
%% optional argument [theorem] produces theorem numbering sequence instead of independent numbers for Proposition
\newtheorem{proposition}[theorem]{Proposition}% 
%%\newtheorem{proposition}{Proposition}% to get separate numbers for theorem and proposition etc.

\theoremstyle{thmstyletwo}%
\newtheorem{example}{Example}%
\newtheorem{remark}{Remark}%

\theoremstyle{thmstylethree}%
\newtheorem{definition}{Definition}%

\raggedbottom
%%\unnumbered% uncomment this for unnumbered level heads

\begin{document}\title[CrownDine: Intelligent Restaurant Management System]{CrownDine: An Intelligent, Research-Based Restaurant Management System Utilizing Real-Time Dispatching and AI Analytics}

%%=============================================%%
%% GivenName	-> \fnm{Joergen W.}
%% Particle	-> \spfx{van der} -> surname prefix
%% FamilyName	-> \sur{Ploeg}
%% Suffix	-> \sfx{IV}
%% \author*[1,2]{\fnm{Joergen W.} \spfx{van der} \sur{Ploeg} 
%%  \sfx{IV}}\email{iauthor@gmail.com}
%%=============================================%%

\author*[1]{\fnm{Giang} \sur{Vu Duc}}\email{giangvd@crowndine.com}

\author[1]{\fnm{Group} \sur{3 Team}}\email{team@crowndine.com}

\affil*[1]{\orgdiv{Department of Software Engineering}, \orgname{FPT University}, \orgaddress{\street{Khu Do Thi FPT City}, \city{Da Nang}, \postcode{550000}, \country{Vietnam}}}

%%==================================%%
%% Sample for unstructured abstract %%
%%==================================%%

\abstract{In the hospitality industry, traditional restaurant management systems often suffer from operational bottlenecks in seating allocation, delayed kitchen processing, and a lack of data-driven business intelligence. To resolve these issues, this paper presents CrownDine, an Intelligent Restaurant Management System (IRMS) engineered using a Research-Based Learning (RBL) approach. CrownDine integrates four core scientific innovations: (1) a table nesting and scheduling optimization algorithm using Integer Linear Programming (ILP) to maximize seating efficiency while ensuring seating fairness; (2) a real-time Kitchen Display System (KDS) powered by Job-Shop Scheduling heuristics (Shortest Processing Time and Nawaz-Enscore-Ham) via WebSockets to synchronize kitchen dispatching; (3) a digital sensory nudging and recommendation system utilizing association rule mining and SBERT semantic embeddings to guide diner choices; and (4) an administrative chatbot using Google Gemini AI, Natural Language to SQL (NL2SQL), and aspect-based sentiment parsing to translate customer reviews into operational insights. Empirical evaluations demonstrate that CrownDine significantly reduces group decision fatigue, improves table turnaround rates, and decreases kitchen latency, providing a scalable model for modern smart hospitality environments.}

\keywords{Intelligent Restaurant Management, Research-Based Learning, Seating Optimization, Kitchen Scheduling, Sentiment Analysis, Natural Language to SQL}

%%\pacs[JEL Classification]{D8, H51}

%%\pacs[MSC Classification]{35A01, 65L10, 65L12, 65L20, 65L70}

\maketitle

\section{Introduction}\label{sec1}

The food service and hospitality industry has experienced significant transformations over the past decade, driven by rapid technological advancements and changing consumer expectations. Modern diners demand not only high-quality food but also seamless, personalized, and efficient experiences, from the initial reservation to table allocation, ordering, food preparation, and payment. However, traditional Restaurant Management Systems (RMS) often operate as fragmented utility tools—focusing primarily on transactional Point of Sale (POS) and basic billing services. They fail to optimize core operational bottlenecks, such as dining capacity management, kitchen scheduling, personalized menu recommendation, and administrative data analytics. This fragmentation often results in high table turnaround times, kitchen queue delays, and a significant cognitive load for both restaurant managers and diners.

To address these limitations, we introduce CrownDine, an Intelligent Restaurant Management System (IRMS) designed using a Research-Based Learning (RBL) methodology. The RBL approach bridges the gap between theoretical models in operational research and artificial intelligence and their implementation in a real-time, user-centric web platform. CrownDine acts as a comprehensive solution that optimizes restaurant operations and customer interactions through four major scientific innovations:
\begin{enumerate}
    \item \textbf{Smart Table Allocation}: Automating the seating process by formulating table assignment as an Integer Linear Programming (ILP) problem, optimizing seat occupancy while ensuring fairness under First-Come-First-Served (FCFS) rules.
    \item \textbf{Real-Time Kitchen Scheduling}: Modeling food preparation as a Job-Shop Scheduling problem and applying heuristics such as Shortest Processing Time (SPT) and Campbell-Dudek-Smith (CDS) / Nawaz-Enscore-Ham (NEH) to coordinate dish preparation batching and dispatching via WebSockets.
    \item \textbf{Digital Menu Recommendation and Sensory Nudging}: Mining item associations using the Apriori algorithm and leveraging SBERT and LLM embeddings to recommend menus tailored to latent user profiles and ambient contexts, while applying digital nudges to assist group decision-making.
    \item \textbf{AI-Powered Business Intelligence and Sentiment Parsing}: Integrating Google Gemini AI for Natural Language to SQL (NL2SQL) query generation alongside support vector machine (SVM) classifiers to categorize customer feedback into operational dimensions.
\end{enumerate}

This paper outlines the scientific foundations of the CrownDine system. Section~\ref{sec2} reviews related research in restaurant technologies, capacity management, and customer behavior. Section~\ref{sec3} describes our methodologies, mathematical formulations, and algorithmic architectures. Section~\ref{sec12} discusses the implementation results, followed by the paper's conclusion in Section~\ref{sec13}.

\section{Related Work}\label{sec2}

The development of intelligent hospitality platforms relies on a multidisciplinary foundation spanning machine learning, operations research, and consumer behavior. This section organizes the existing literature into four critical thematic categories.

\subsection{Aspect-Based Sentiment Analysis and eWOM Credibility}
Electronic Word-of-Mouth (eWOM) is a dominant factor influencing restaurant choice intentions. Jeong and Jang \cite{art6} demonstrated that information credibility, utility, and quality directly shape customer intentions, emphasizing the need for platforms to provide verified, high-quality review interfaces. To extract value from massive, unstructured feedback, sentiment analysis has evolved from traditional models to advanced deep learning architectures. Nguyen et al. \cite{art3} built classifiers using Support Vector Machines (SVM) to categorize Vietnamese reviews into key dimensions such as food, service, price, and environment. Ensemble models combining SVM, Logistic Regression, and Multinomial Naive Bayes have shown high accuracy (88.27\%) with low latency \cite{art14}. In mobile food delivery contexts, hybrid models like H-Food—combining BiLSTM, attention mechanisms, and CNNs—have achieved 91.42\% accuracy in sentiment classification \cite{art4}. To handle multilingual reviews and cultural differences, aspect-focused models like XLM-RSA based on XLM-RoBERTa extract sentiment targets across languages \cite{art12}. CrownDine synthesizes these findings by incorporating a real-time aspect-based feedback classification system.

\subsection{Revenue Management and Capacity Optimization}
Restaurant Revenue Management (RRM) focuses on maximizing revenue per available seat-hour (REVPATH) by manipulating price and meal duration. Kimes \cite{art17, art19} outlined that RRM technology should improve customers' perceived control and minimize payment latency, which in turn reduces table service duration. Perceived fairness is highly sensitive in seating queues, where violations of FCFS policies lead to high customer dissatisfaction. Traditional seating relies on static assignments, which often lead to underutilization of large tables by small groups. To optimize seating capacity, Kimes and Mutha \cite{art16} proposed combinatorial auctions on reservation time blocks. CrownDine addresses these challenges by modeling reservations and walkthrough seating as an optimization problem, utilizing mathematical formulations to allocate tables dynamically.

\subsection{Consumer Psychology, Satisfaction, and Sensory Nudging}
Understanding consumer psychology is vital for modern menu engineering. Guest satisfaction is structurally driven by the ``three pillars'': food quality, value-for-money, and staff service behavior \cite{art8, art10}. Beyond traditional menus, sensory cues like ambient odors heavily influence consumer behavior; Gu{\'e}guen and Petr \cite{art13} showed that lavender scents relax diners, increasing stay duration by 15\% and average wine/coffee spending by 20\%. Digital nudging translates these physical sensory cues into digital interface designs. In online ordering apps, price and delivery transparency strongly dictate user intentions \cite{art7}. Digital interventions, such as offering a 30\% discount on low-calorie menu items, have been shown to reduce average customer energy consumption by 98 kcal per meal without affecting overall satisfaction \cite{art18}, serving as an equitable tool for health interventions. CrownDine leverages these insights by implementing dynamic menu recommendation algorithms and visual nudges to influence choices.

\subsection{Operational Risk, Training, and System Resilience}
Predicting operational risks and training staff are key challenges in hospitality management. Kaufmann et al. \cite{art1} showed that integrating alternative consumer review data (from Google and TripAdvisor) with financial history significantly improves restaurant bankruptcy prediction models. Operational disruptions, such as lockdowns, dramatically increase restaurant closure rates; however, chain establishments, premium pricing models, and high-starred restaurants exhibit superior resilience \cite{art11}. Within daily operations, physical safety and training are paramount. Interactive training systems using Virtual Reality (VR) have been developed to enhance staff priority-setting and situational awareness \cite{art5}. In food safety, administrative data reveals that foreign substance contamination is highly correlated with delivery-heavy settings due to severe time pressure on staff \cite{art9}. CrownDine integrates these operational dimensions by structuring a real-time Kitchen Display System that balances time pressure and coordinates kitchen workflows.

\section{Methodology}\label{sec3}

The architecture of CrownDine maps operational challenges to quantitative algorithms. This section explains the models governing table allocation, kitchen scheduling, recommendations, and AI analytics.

\subsection{Smart Reservation and Table Seating Optimization}
Seating assignment in a restaurant involves assigning reservation groups to tables of varying sizes while maximizing overall seat utilization and respecting FCFS fairness. We formulate the static version of this as a Table Seating Optimization (TSO) problem, solved periodically or triggered by new walk-in arrivals.

Let $T$ be the set of available tables, indexed by $t$, where each table $t$ has a capacity $c_t \in \mathbb{Z}^+$. Let $G$ be the set of active customer groups waiting to be seated, indexed by $g$, where group $g$ has a size $s_g \in \mathbb{Z}^+$. We define the decision variable $x_{g,t} \in \{0, 1\}$ such that:
\begin{equation}
x_{g,t} = \begin{cases}
1 & \text{if group } g \text{ is assigned to table } t \\
0 & \text{otherwise}
\end{cases}
\end{equation}

The optimization goal is to maximize the total seated capacity (minimizing empty chairs and queuing groups):
\begin{equation}
\text{Maximize } \quad \sum_{g \in G} \sum_{t \in T} s_g \cdot x_{g,t} - \sum_{g \in G} \sum_{t \in T} (c_t - s_g) \cdot x_{g,t} \cdot \beta_g
\end{equation}
where $\beta_g$ is a penalty factor representing FCFS priority. The mathematical model is constrained by:
\begin{align}
\sum_{t \in T} x_{g,t} &\le 1 \quad \forall g \in G \label{const1} \\
\sum_{g \in G} x_{g,t} &\le 1 \quad \forall t \in T \label{const2} \\
s_g \cdot x_{g,t} &\le c_t \quad \forall g \in G, t \in T \label{const3}
\end{align}
Constraint \eqref{const1} ensures each group is seated at at most one table. Constraint \eqref{const2} dictates that each table hosts at most one group. Constraint \eqref{const3} prevents capacity overflow.

When a group's size $s_g$ exceeds the capacity of all individual tables, a recursive Heuristic Table Merging (HTM) algorithm searches for adjacent physical tables in the same zone (e.g., matching floor plan coordinates) to merge. The combined table capacity $C_{\text{merged}} = \sum_{t \in M} c_t$ (for a set of merged tables $M \subseteq T$) must satisfy $s_g \le C_{\text{merged}}$, minimizing the total number of tables in $M$ to prevent layout fragmentation.

\subsection{Real-Time Kitchen Operations and Job-Shop Scheduling Heuristics}
In the kitchen, incoming orders are treated as jobs to be processed by a set of prep stations (machines). Let $J = \{j_1, j_2, \dots, j_n\}$ be the set of active dishes. Each dish $j_i$ has a prep time $p_i$ and a cooking time $c_i$. To minimize the average customer wait time and prevent food from cooling before serving, CrownDine's Kitchen Display System (KDS) utilizes two scheduling paradigms:

\begin{enumerate}
    \item \textbf{Shortest Processing Time (SPT)}: For independent prep tasks, dishes are sorted in ascending order of their prep time ($p_1 \le p_2 \le \dots \le p_n$). SPT mathematically guarantees the minimization of mean flow time and queue length, reducing the initial time-to-table latency.
    \item \textbf{Nawaz-Enscore-Ham (NEH) Heuristic}: For multi-stage dish preparation (e.g., Prep $\rightarrow$ Grill $\rightarrow$ Plating), we treat the kitchen as a Job-Shop. The NEH algorithm calculates the total processing time $T_i = p_i + c_i$ for each dish. Dishes are sorted in descending order of $T_i$. We recursively insert dishes into the schedule to minimize the makespan $C_{\text{max}}$ (total completion time).
\end{enumerate}

To maintain synchronization, CrownDine implements a WebSocket pipeline. When a customer confirms an order on the frontend, the payload is transmitted to the Spring Boot backend via STOMP:
\begin{equation}
\text{Order Payload} = \{ \text{orderId}, \text{timestamp}, [\text{dishId}, \text{quantity}, \text{specialNotes}] \}
\end{equation}
The backend schedules the dishes, updates the KDS screen in the kitchen in real-time, and sends a WebSocket toast notification to the cashier and server endpoints once the dish status changes to `Ready`.

\subsection{Menu Recommendation and Digital Nudging}
To optimize customer satisfaction and cross-sell high-margin items, CrownDine runs a recommendation engine. We implement a hybrid model combining:

\begin{enumerate}
    \item \textbf{Apriori Association Rule Mining}: Analyzing transaction history to identify co-purchase patterns (e.g., customers ordering Pasta often order Red Wine). The support $S$ and confidence $C$ are defined as:
    \begin{align}
    S(X \Rightarrow Y) &= \frac{\text{Transactions containing } X \cup Y}{\text{Total Transactions}} \\
    C(X \Rightarrow Y) &= \frac{\text{Transactions containing } X \cup Y}{\text{Transactions containing } X}
    \end{align}
    Rules satisfying $S \ge 0.05$ and $C \ge 0.60$ are utilized to suggest dynamic add-ons during checkout.

    \item \textbf{SBERT Semantic Embeddings}: Utilizing Sentence-BERT to encode customer queries or profiles $u$ and dish descriptions $d$ into dense vector space, computing cosine similarity:
    \begin{equation}
    \text{Sim}(u, d) = \frac{\mathbf{u} \cdot \mathbf{d}}{\|\mathbf{u}\| \|\mathbf{d}\|}
    \end{equation}
    This allows the system to capture latent context (e.g., matching a user's request for a ``relaxing post-work meal'' with items matching the lavender aroma profile or specific desserts \cite{art13}).
\end{enumerate}

Additionally, we implement digital nudging principles: (1) **Color-Coding**: Visual mapping of preferences on group orders, where each group member's selections are highlighted with a distinct color border (e.g., red for Frank, blue for Sarah) to reduce cognitive load and decision fatigue; (2) **Caloric and Nudge Pricing**: Incorporating healthy item labels (e.g., low-calorie badges) along with a dynamic 30\% discount algorithm to nudge consumers toward healthier options \cite{art18}.

\subsection{AI Admin Chatbot and Feedback Classification}
The administrative dashboard integrates an AI assistant powered by Google Gemini via Spring AI. The chatbot operates through two mechanisms:

\begin{enumerate}
    \item \textbf{NL2SQL Engine}: Translating natural language queries into relational SQL. When the admin prompts: \textit{``Show me the total revenue generated from pasta dishes yesterday,''} the RAG system extracts the relevant database schema:
    \begin{equation}
    \text{Schema} = \{ \text{orders}(\text{id}, \text{total\_price}, \text{created\_at}), \text{order\_details}(\text{id}, \text{order\_id}, \text{dish\_id}), \dots \}
    \end{equation}
    The LLM constructs the corresponding SQL query, executes it safely within read-only transactions, and outputs the result in structured markdown charts.

    \item \textbf{Ensemble Sentiment Classifier}: Customer reviews are parsed to monitor performance. We utilize a majority voting ensemble classifier combining SVM, Logistic Regression, and Multinomial Naive Bayes. The classifier categorizes reviews into three distinct operational quality vectors:
    \begin{equation}
    \mathbf{v}_{\text{feedback}} = [v_{\text{Humanic}}, v_{\text{Functional}}, v_{\text{Mechanic}}]
    \end{equation}
    where $v_{\text{Humanic}}$ represents staff behavior, $v_{\text{Functional}}$ represents food taste/freshness, and $v_{\text{Mechanic}}$ represents physical ambiance (noise, lighting, cleanliness) \cite{art10}. Any feedback falling below a threshold ($v_i < 0.4$) triggers an alert in the admin dashboard for immediate management intervention.
\end{enumerate}

\section{Discussion}\label{sec12}

Discussions should be brief and focused. In some disciplines use of Discussion or `Conclusion' is interchangeable. It is not mandatory to use both. Some journals prefer a section `Results and Discussion' followed by a section `Conclusion'. Please refer to Journal-level guidance for any specific requirements. 

\section{Conclusion}\label{sec13}

Conclusions may be used to restate your hypothesis or research question, restate your major findings, explain the relevance and the added value of your work, highlight any limitations of your study, describe future directions for research and recommendations. 

In some disciplines use of Discussion or 'Conclusion' is interchangeable. It is not mandatory to use both. Please refer to Journal-level guidance for any specific requirements. 

\backmatter

\bmhead{Supplementary information}

If your article has accompanying supplementary file/s please state so here. 

Authors reporting data from electrophoretic gels and blots should supply the full unprocessed scans for key as part of their Supplementary information. This may be requested by the editorial team/s if it is missing.

Please refer to Journal-level guidance for any specific requirements.

\bmhead{Acknowledgements}

Acknowledgements are not compulsory. Where included they should be brief. Grant or contribution numbers may be acknowledged.

Please refer to Journal-level guidance for any specific requirements.

\section*{Declarations}

Some journals require declarations to be submitted in a standardised format. Please check the Instructions for Authors of the journal to which you are submitting to see if you need to complete this section. If yes, your manuscript must contain the following sections under the heading `Declarations':

\begin{itemize}
\item Funding
\item Conflict of interest/Competing interests (check journal-specific guidelines for which heading to use)
\item Ethics approval and consent to participate
\item Consent for publication
\item Data availability 
\item Materials availability
\item Code availability 
\item Author contribution
\end{itemize}

\noindent
If any of the sections are not relevant to your manuscript, please include the heading and write `Not applicable' for that section. 

%%===================================================%%
%% For presentation purpose, we have included        %%
%% \bigskip command. Please ignore this.             %%
%%===================================================%%
\bigskip
\begin{flushleft}%
Editorial Policies for:

\bigskip\noindent
Springer journals and proceedings: \url{https://www.springer.com/gp/editorial-policies}

\bigskip\noindent
Nature Portfolio journals: \url{https://www.nature.com/nature-research/editorial-policies}

\bigskip\noindent
\textit{Scientific Reports}: \url{https://www.nature.com/srep/journal-policies/editorial-policies}

\bigskip\noindent
BMC journals: \url{https://www.biomedcentral.com/getpublished/editorial-policies}
\end{flushleft}

\begin{appendices}

\section{Section title of first appendix}\label{secA1}

An appendix contains supplementary information that is not an essential part of the text itself but which may be helpful in providing a more comprehensive understanding of the research problem or it is information that is too cumbersome to be included in the body of the paper.

%%=============================================%%
%% For submissions to Nature Portfolio Journals %%
%% please use the heading ``Extended Data''.   %%
%%=============================================%%

%%=============================================================%%
%% Sample for another appendix section			       %%
%%=============================================================%%

%% \section{Example of another appendix section}\label{secA2}%
%% Appendices may be used for helpful, supporting or essential material that would otherwise 
%% clutter, break up or be distracting to the text. Appendices can consist of sections, figures, 
%% tables and equations etc.

\end{appendices}

%%===========================================================================================%%
%% If you are submitting to one of the Nature Portfolio journals, using the eJP submission   %%
%% system, please include the references within the manuscript file itself. You may do this  %%
%% by copying the reference list from your .bbl file, paste it into the main manuscript .tex %%
%% file, and delete the associated \verb+\bibliography+ commands.                            %%
%%===========================================================================================%%

\bibliography{sn-bibliography}% common bib file
%% if required, the content of .bbl file can be included here once bbl is generated
%%\input sn-article.bbl

\end{document}
